\documentclass{article}
\usepackage{amsmath}
\title{Retained closure, a consuming cut, and reassignment of a backup window}
\date{}
\begin{document}
\maketitle

\paragraph{Scope and source correspondence.}
An old delivery stage has a published identifier and a reserved way
to retry. A source must remain out of a named bay until it receives
the qualified terminal result. Another job could use the reserved
radio transmitter before the retry deadline. The finite reference
joins actual RWARE source and station transitions, the existing
source-held authorization rule, a transactional station worker, and
pinned SubT dispatch and rolling radio accounting. Versions and
selected executable fragments are pinned in \verb|SOURCE.tsv|. All
effects, debits and delivered evidence in each result are on the
same run. The store is authoritative over a simulator transition,
not over external motor motion. Application packet names, the link
availability seam, endpoint topology, clock relation, receipt
semantics and workload are supplied by this development adapter.
The controller cannot inspect the station store or a fault parameter:
the station makes its own authoritative receipt and the controller
observes actual callbacks and local debit.

\paragraph{Fixed positive operating envelope.}
The mandatory floor is the old station result by tick seven and
local release by tick eight, with two independent motions and
radio callbacks preserved. At tick two the old request reaches the
worker, which exits before its native step, after the private step,
after staging the durable transaction, or after the commit. The
caller receives no direct result in all four cases. Q's actual
setup arrives at four and provides prepared authority at six; its
station route and the final terminal path are available in the
positive envelope. P's station route can drop at seven, hence
before any published effect in the first three crash branches.
The same common interruption may also lose an earlier Base probe
at three. That send spends four bytes at Base, even without a
callback. No optional effect is part of the mandatory floor.
Setup and backup route loss are tested separately but not assumed
to preserve the floor.

\paragraph{Composing the decisions.}
A first receipt attempt at six uses the actual canonical durable
record. It offers a matching result when the old effect is
settled; otherwise a two-byte older acknowledgement for a committed
predecessor can arrive. Under the common interruption there is
only one byte of Base headroom at six, so either result is masked.
Thus identical delivered setup and silence have two concrete
explanations even under the same interruption: an old effect
still pending, or already durable. Restart safety comes from using
the same key and the atomic store in both histories. It does not
give permission to infer which history occurred.

A three-byte optional send on Q and a three-byte backup retry both
fit the five-byte window separately. Without a matching receipt,
the retry may be necessary, and both actions cannot fit. A genuine
result discharges that particular debt and permits the optional
send on Q. It also ends the bay hold: the original carrier actually
returns before the optional callback. The old clearance cannot then
be reused for a second carrier. In the negative control the latter
carrier requests its stock forward move and is blocked, even though
the communication callback succeeds. In another negative control
an unrelated predecessor receipt releases the optional radio
capacity early and can leave the old effect unable to meet its
deadline. The arithmetic rule, including the incompatibility of
the two locally valid checks, is verified in \verb|Joint_Admission.lean|.
It assumes the named matching condition and live-window accounting;
Lean is not used to claim verification of either simulator.

\paragraph{Constructive service argument.}
If the matching result arrives at six, the effect is already
published and the bay is released only after that callback. If
not, the source stays out and no optional debit is placed on Q.
An available P completes or returns the cached result at seven.
If P's path drops, Q is ready, Q's transmitter still has all five
bytes, and Base's old incoming request is older than one radio
second. Q therefore reaches the effect handler with the original
key, committing once or returning the original result. At eight
the four-byte failed debit from three is outside the inclusive
window. A successful nonmatching two-byte predecessor receipt
at six fits with the new three-byte terminal result in both
Base TX and source RX. The authentic current result then permits
source return. All three positive paths preserve the independent
motion and traffic, which use separate cells and receive budgets.
A much later duplicate retry is executed to check that the
committed native state and credit remain exactly one step.

\paragraph{Comparison boundary.}
The unchanged shared core and the direct rule with the same
information and resources agree trace for trace. Each of the four
arms executes 64 cases, 16 in the declared positive envelope. Both
typed controllers meet the full floor in all 16 and settle and return
on time in 37 of the full 64, with eight lawful optional sends. The
core requests 1,419 radio bytes and actually debits 1,347, including
the scripted follow-up retries. The receiver sees 37 on-time returns,
not a probability estimate. The untyped arm misses three of the
16 floor cases; the stale-clearance arm produces eight inaccurate
admission claims while the native motion remains blocked. The untyped and
stale-clearance arms are diagnostics, not claims of an improvement
over a strong published assembly. The gate treats the retained
old state, source motion and relayed settlement as different
obligations. No independent transfer or inventory-wide coverage
follows from this reference family.
\end{document}
